Corvus is run from within a terminal (or miniforge dos command window on windows). If you used the 
install script to setup corvus, there should be an executable script named "Corvus-version.command" or "Corvus-version.bat" 
on your Desktop. To run corvus, double click this scipt, navigate to the input file you want 
to run, and type,
\begin{verbatim}
run-corvus -i your_input_file
\end{verbatim}

If you installed with the step-by-step instructions, you will need to activate your
corvus environment first. With conda, this is done by typing,
\begin{verbatim}
conda activate your_corvus_env
\end{verbatim}

\subsection{Example files}
There are a variety of example input files located in the folder
\begin{verbatim}
your_home_directory/corvus_examples
\end{verbatim}
If you installed with the install.sh or install.bat scripts, your terminal
will start in the examples folder when you use the launcher provided on 
your desktop. You can then run a set of the input files using the run_examples.py
python script,
\begin{verbatim}
python run_examples.py
\end{verbatim}

You can also compare various files with their references, which have
the same name with REFERENCE\_ prepended to it, e.g., 
REFERENCE_Corvus.xanes.out
